\begin{abstract}
  
While softmax attention is widely used recently, its quadratic complexity leads
to high computational costs. Recent linear attention mechanisms mitigate the quadratic bottleneck but still redundantly process uninformative background tokens, wasting computational budget and limiting the model accuracy. In this study, we propose a gating framework to evaluate the importance of each token. Specifically, we use Multi-Agent Reinforcement Learning (MARL), where each visual token acts as an autonomous agent that utilizes a shared Actor-Critic policy to output continuous survival weights, gating information flow into the recurrent memory. For image classification tasks, our model reached accuracies of \(48.88 \% \) and \(58.37 \% \) on the datasets ImageNet100 and CIFAR-10 respectively. While the framework remains in its early stages of development and current performance metrics are modest, the resulting attention maps suggest that this gating approach holds potential for future optimization.
  % Use the word ``Abstract'' as the title, in 12-point Times, boldface type, centered relative to the column, initially capitalized.
  % The abstract is to be in 10-point, single-spaced type.
  % Leave two blank lines after the Abstract, then begin the main text.
  % Look at previous \confName abstracts to get a feel for style and length.
\end{abstract}